\clearpage
\section{Thirty-six standalone attack questions}
\label{sec:questions}

Each setup describes the system needed for the experiment. These are questions for public
interface knowledge or a self-contained reproducer, not instructions to inspect a private
trace or a repository-specific state machine. Proposed observations are not claimed results.

\question{Must a small edit require a whole-space scan?}
{An emulator cannot intercept stores to guest GPU page tables. After a translation barrier it must
discover changes before dependent commands run, although only a small range may have changed.}
{Identify a public changed-range certificate or bounded validation rule, then try a root or
subtree change outside its claimed scope.}
{A certificate that survives that counterexample permits incremental work. Otherwise the
whole-space scan remains part of this design's correctness cost.}
\label{q:whole}

\question{Which acknowledgement makes old backing reusable?}
{GPU work may still refer to a VA while software removes that host mapping, invalidates
translation caches and wants to reuse its physical backing.}
{State the public host-driver ordering guarantee, or run a counter test that observes an old
access after the proposed retirement event.}
{A weaker guarantee requires a separate execution/lifetime fence; a sufficient one can simplify
retirement without relying on elapsed time.}
\label{q:retirement}

\question{Does a longer observation window show late work or replacement work?}
{A fixed Windows graphics workload is run with default recovery settings and a longer diagnostic
window. Host submissions and guest recovery each carry a lifetime identity.}
{Determine whether the original operation eventually completes, is canceled, or is replaced by new
work after recovery. Locate the scheduler preemption request too.}
{Only identified original completion supports a lateness explanation. Replacement or missing
identity keeps the lost-completion explanation open.}
\label{q:watchdog}

\question{Can live sparse reservations be combined without a hole?}
{Two adjacent GPU VA reservations intentionally absorb unbacked accesses. Software wants one
larger reservation without changing any currently accessible address or pausing execution.}
{Find a public atomic replacement operation, or demonstrate a transition that avoids
release-before-reserve faults under concurrent accesses.}
{A valid primitive permits live geometry compression. Without it, empty-reservation recycling does
not justify merging active sparse coverage.}
\label{q:coalesce}

\question{What bounds a single host-driver call?}
{A software budget can stop starting operations, but cannot interrupt an already running native
memory-reservation, map or invalidate call.}
{Measure per-call tails under idle and loaded conditions and identify any public bound,
cancellation mechanism or asynchronous alternative.}
{A usable bound enables deadline reasoning. Otherwise dependent work must tolerate an unbounded
tail or the dependency must be moved ahead of the guest-visible barrier.}
\label{q:nativecall}

\question{What must a lock holder promise before a worker parks?}
{A worker consumes a notification, fails a try-lock, and prepares to sleep while work remains. The
holder may release before, during or after waiter registration.}
{Provide a register/recheck/release protocol and enumerate all three orderings, including
retirement and reuse of the protected object.}
{A missed final wake changes the handoff protocol. A correct protocol permits sleep without using
a periodic timer as the actual source of progress.}
\label{q:wake}

\question{What separates GPU clients running at different Unix UIDs?}
{Two VMM processes have distinct unprivileged Unix identities, each holding a host GPU driver
client and USER-privilege channels in separate VA spaces.}
{Document or test which cross-client operations remain possible without explicit resource sharing,
including direct naming and resource duplication.}
{This determines what the credential policy actually closes and which operations still need
application-level ownership checks or a stricter deployment boundary.}
\label{q:uid}

\question{How much of a nested result belongs to nesting?}
{The same GPU workload can run natively, in a guest on a physical host, or in a guest whose host
is itself a VM. Submission count, placement and driver versions can be held fixed.}
{Measure exit counts, per-exit time, computation time and output in all three configurations while
recording clock state and host load.}
{The decomposition tells which overhead is an emulator mechanism and which ratio cannot be
transferred from a rented nested host to a physical desktop.}
\label{q:nested}

\question{Which oracle distinguishes correct GPU execution from a plausible PASS?}
{A harness receives an application exit code, output bytes and an activity signal. A process can
exit zero without checking a failed GPU synchronisation.}
{Inject separately a missing launch, wrong output, forced GPU error and CPU fallback; determine
which independent observation rejects each while accepting the native control.}
{This selects the evidence needed for the workload claim, rather than assuming either an activity
counter or output equality alone proves the intended execution path.}
\label{q:proof}

\question{What can a rented-machine GPU gate establish?}
{A provider operator can inspect approved source and possesses one run-scoped reporting
capability. It can fabricate logs and a correctly bound structured result.}
{Separate errors a normal run can expose from claims requiring an observer or builder outside that
operator's control.}
{The answer defines the useful regression tripwire and prevents it from being treated as release
provenance or an execution attestation.}
\label{q:rental}

\question{What happens to the host desktop when a guest fails?}
{A host display session and a guest use one physical GPU. The host retains a live rendering
context and reference data while the guest faults, hangs work or triggers recovery.}
{Measure host frame latency, output integrity, context survival and recovery; compare the same
disturbance generated by an ordinary native host client.}
{A host-wide failure restricts sharing claims even if guest addresses are isolated. Containment
supports only the exercised fault classes and measured service bounds.}
\label{q:host}

\question{Could a mapping barrier acknowledge before host apply?}
{A guest edits GPU page tables, asks for an invalidate, and expects subsequent work to see the
edit. The emulator still has an old host mapping and no universal fault/replay service.}
{Construct either a safe early-ack protocol or a stale-access counterexample, distinguishing an
absent mapping that faults from a still-valid old mapping that returns wrong data.}
{This decides whether blocking apply is inherent in the current authority model, or whether a
separately proven fault/retry or versioning mechanism could remove it.}
\label{q:early}

\question{How should driver-release maintenance cost be measured?}
{An emulator serves an unmodified guest driver while independently selecting host-driver ABIs and
die capabilities. New releases can change layouts, requests and behaviour.}
{For a sequence of public driver releases, record changed interfaces, implementation and test
effort, unavailable hardware cells and remaining hand-maintained answers.}
{The resulting ledger can price a supported release window and proactive measurements; a
compatibility matrix alone cannot estimate that maintenance burden.}
\label{q:releasecost}

\question{Which software agreements apply to this deployment?}
{A source-built emulator runs a vendor driver on GeForce-class hardware, possibly on rented
infrastructure or as a commercial service; its own source has separate licence exceptions.}
{Obtain qualified review of the actual driver agreement, platform/hosting/datacentre terms,
technical restrictions and combined-binary distribution permissions for that arrangement.}
{The review determines release and deployment scope. Neither a working VM nor absence of a vendor
vGPU component answers the contractual questions.}
\label{q:eula}

\question{What is the worst barrier cost the guest can cause?}
{One translation barrier triggers a table traversal plus host maps, reservation changes, repair
and invalidation. A guest can choose a fragmented layout within finite object caps.}
{Bound or measure each non-overlapping stage as a function of visited tables and native
operations, then identify every host or guest client waiting behind the operation.}
{This decides whether workload caps imply a useful latency/service bound or merely finite memory
use with a potentially unacceptable shared delay.}
\label{q:worst}

\question{Could isolated hardware translation remove the software walk?}
{Instead of reconstructing host GPU tables, imagine a host interface that accepts a guest
translation root in a hardware-isolated context. The guest can mutate or recycle it.}
{Identify a public interface and prove ownership, privilege, revocation and concurrent-update
semantics, including all privileged addressing paths.}
{A usable interface changes the architecture. A root-validation check without ongoing isolation is
not enough to replace mirroring safely.}
\label{q:adopt}

\question{What can a walk promise while tables change concurrently?}
{A GPU traversal reads a directory, then the guest replaces or recycles a lower table before the
traversal reaches it. Unchanged addresses are supposed to remain usable.}
{Give a protocol and counterexamples for allowed old/new/fault outcomes, root generation reuse,
partial reads and a changed subtree shared by two roots.}
{This determines whether live walking needs validation, versioning or a stronger contract;
treating every transient decode failure as a stable malformed mapping is insufficient.}
\label{q:race}

\question{How much coupling is inside the shared host driver?}
{Two clients use independent GPU VA spaces but issue memory-management calls through one host
driver. One client repeatedly creates difficult reservation layouts.}
{Measure the other client's call latency and host-display latency, distinguishing per-client
queues from any driver-wide serialisation and GPU execution contention.}
{This determines whether per-space emulator gates isolate useful latency or leave a shared native
bottleneck requiring admission or deployment limits.}
\label{q:global}

\question{Which limits give a VM a useful service guarantee?}
{A guest has a fixed video-memory allocation but shares GPU scheduling, encoder capacity and
host-driver request service with a live host workload.}
{Vary allocations, scheduling requests and request rate independently; measure the victim's worst
service gap and whether releasing the load restores service without a reset.}
{This separates memory isolation from fairness and identifies the quota or admission contract
required for a stated responsiveness target.}
\label{q:quota}

\question{What state must a GPU VM snapshot or migration contain?}
{A VM is stopped with GPU work in flight, host-owned memory objects and queued completions. It may
resume in another process or on a different compatible GPU.}
{Specify which state can be drained, reconstructed or transferred and how old completions, host
identities and vendor-driver state are prevented from crossing lifetimes.}
{This identifies a supportable snapshot/migration mode or makes its exclusion explicit, rather
than assuming CPU-memory snapshot semantics cover the device.}
\label{q:snapshot}

\question{Can unaffected ranges submit through a refresh?}
{One host VA space runs work on range A while the guest changes disjoint range B. A global gate
currently withholds all new submissions into that space.}
{Provide a bounded range-dependency certificate that survives root changes, aliases, publication
races and a full change-history buffer.}
{It permits narrower admission if sound; otherwise the implementation must keep conservative holds
and reduce the work behind them.}
\label{q:range}

\question{What preserves a running command ring across relocation?}
{A stock driver requests a new GPU command-reference ring while commands from the old ring are in
flight. Software wants native relocation rather than relaying every command.}
{Demonstrate cursor continuity, ownership of both ring memories, retirement of the old source and
a completion produced only by the replacement ring.}
{This selects between a supported native move, software mediation from birth and a separate
cooperation tier without relying on an idle cursor as completion proof.}
\label{q:relocate}

\question{Which clock separates GPU completion from CPU observation?}
{Software timestamps a submission and its first read of a completed GPU-written marker. The
observing worker may spend time in a native driver call.}
{Find an independent GPU-time witness and clock relation, or a second observer that distinguishes
execution delay from delayed sampling without adding a correctness poll.}
{This determines which latency component can be reported and which must remain an ambiguous
submission-to-observation sum.}
\label{q:clock}

\question{How can a notification be tied to one operation?}
{Several channels and clients share engine-wide completion notifications. A CPU may see an
interrupt before it samples its own completion marker.}
{Specify the minimum identity, masking/coalescing state and ordering evidence needed to establish
that a notification covered the particular operation.}
{This chooses between a valid delivery-latency sample and an explicitly incomplete correlation; an
unrelated interrupt must not manufacture a completed record.}
\label{q:interrupt}

\question{What happens when a finite sparse budget is exhausted?}
{A guest intentionally maps unbacked zero/absorbing GPU pages. The host has a finite reservation
budget and leaves additional requested coverage invalid on exhaustion.}
{Identify a recoverable driver-visible refusal point, or a workload that requires those semantics
after that point and cannot recover from an invalid mapping.}
{This defines whether the bounded mode is compatible, needs a new representation or must exclude
that workload; counting the deficit alone does not choose among them.}
\label{q:sparse}

\question{Which identity survives root replacement and reset?}
{A guest VA-space object can replace its page-directory root, reuse an old root address and reset
a channel while host mappings or completions remain outstanding.}
{Give an identity/epoch scheme and adversarial sequence that distinguishes same object, new root
and new lifetime, including delayed callbacks.}
{This determines when a cached mapping, source or completion can be reused and when old
obligations must be retired or canceled.}
\label{q:root}

\question{Which layer fails on resume?}
{A guest resumes from S3 with the emulated GPU state retained, but platform timers and the
operating system also need to recover.}
{Compare matched plain-VMM and device-enabled guests, then require post-wake content-checked GPU
work and locate the first differing state or wait.}
{A platform-only failure changes the harness diagnosis; a device-specific failure remains a device
lifecycle defect even if its firmware-facing state machine is running.}
\label{q:resume}

\question{What distinguishes an ARM port from a cross-build?}
{Logic code builds for arm64, while the host adapter must register shared GPU memory, handle
different page sizes and deliver interrupts and resets.}
{Execute the required memory and lifecycle primitives on a VM-capable ARM host and identify which
ordering or fault-resumption semantics differ from the x86 adapter.}
{This determines whether an adapter capability is missing or whether ordinary porting work is
sufficient. Compilation alone does not choose between them.}
\label{q:arm}

\question{What must a second VMM adapter guarantee?}
{A GPU device core needs bounded traps, shared-memory views, notification and distinct
reset/lifetime events. It currently runs inside one VMM process.}
{Express those guarantees using another public VMM/device protocol, then identify any operation
whose latency or semantics cannot satisfy the same contract.}
{This prices real hypervisor independence and exposes requirements hidden by the first adapter,
without requiring the reader to know its internal traits.}
\label{q:vmm}

\question{How should simultaneous host/guest use be graded?}
{A screenshot shows a host desktop and a guest both using one GPU. The host can also run a
repeatable render or compute workload with an independent output check.}
{Measure each alone and together, preserving clocks, settings and content; report throughput and
tail frame/completion gaps for both sides over repeated boots.}
{This replaces visual coexistence with a bounded shared-use claim and shows whether the advertised
benefit costs unacceptable host responsiveness.}
\label{q:coexist}

\question{Which driver/version cells are most discriminating?}
{Host and guest driver versions are independent. Public releases can change structures inside a
branch and add fallback requests only on error paths.}
{Select cells using actual public interface changes and test that selection against normal and
refusal-driven paths, rather than only major-version labels.}
{This determines the supported window and where a decoder oracle is insufficient without a real
workload and recovery test.}
\label{q:versions}

\question{What is a usable mode for an unmeasured die?}
{An emulated GPU can withhold optional capabilities when their exact-die provenance is
unavailable, while keeping the guest driver unmodified.}
{Identify a minimal public capability surface, exercise useful applications and recovery, and use
wrong-answer mutations to show which advertised fields actually matter.}
{This defines a working reduced mode rather than treating a named refusal or a boot as a complete
compatibility result.}
\label{q:fallback}

\question{Which failures of a viewer may affect a guest?}
{A guest submits display flips while a separate viewer receives copied frames. The viewer can
stall, restart, flood requests or retain its last frame slot indefinitely.}
{Specify and test which guest completions must continue, which frames may be dropped, and the
bounded ownership of buffers across reconnect and reset.}
{This separates guest display semantics from transport quality and determines where a viewer
failure may legitimately cause degraded output rather than a guest hang.}
\label{q:display}

\question{Which independent invariant would expose a shared test-model error?}
{An event-driven queue implementation and its reference model were written from one informal
progress story. Both could forget the same transition.}
{Derive a conservation rule from externally visible submissions, cancellations and effects; plant
a dropped final obligation and show the oracle rejects it.}
{This tests oracle independence through a concrete counterexample, instead of assuming separate
code or multiple reviewers imply separate reasoning.}
\label{q:oracle}

\question{Which repetitions estimate the user-visible failure?}
{Iterations in one boot are stable, but boot state, host placement and concurrent load can change.
Timeouts and failed boots have no final throughput sample.}
{Separate within-boot, between-boot and between-host spread while retaining failures and missing
samples in the denominator.}
{This chooses an acceptance sample unit and prevents a favourable conditional median from
representing a configuration that users cannot reliably reach.}
\label{q:sampling}

\question{What does a GPU USER-privilege test actually prove?}
{A host client can create channels with a hardware privilege level, but a guest also names objects
and supplies addresses through a software decoder.}
{Test physical and virtual operands on each relevant engine and trace the trusted allocation fact
authorising privileged work, independently of guest handles and Unix UID.}
{This establishes the scope of the hardware boundary and the remaining software authority checks;
one accepted ordinary request cannot prove all three permissions.}
\label{q:privilege}

